\chapter{Estado del arte}\label{cap:estado}

En esta sección se va a proceder a realizar un estudio del estado del arte en el ámbito de la planificación de 
recursos empresariales (Enterprise Resource Planning o ERP) para el sector de la oliva.Se analizarán cuáles 
son sus funcionalidades y limitaciones, así como las tecnologías utilizadas para su desarrollo.

\section{Contexto}

El desarrollo de software ERP para el sector de la oliva surge de la necesidad de integrar y controlar procesos 
muy específicos propios de una cadena agroindustrial compleja, caracterizada por una elevada regulación legislativa 
y por unos costes de producción y transformación significativos. La correcta gestión de la información se convierte 
en un factor clave para garantizar la eficiencia operativa, el cumplimiento normativo y la rentabilidad económica de 
las explotaciones, cooperativas y empresas del sector.

En este contexto, España se posiciona como líder mundial en la producción de aceite de oliva, concentrando de media en 
los últimos años más del 40\% de la producción mundial. A gran distancia le siguen otros países productores como Italia, 
con aproximadamente un 9,7\%, Grecia con un 8,2\% y Túnez con un 6,1\%, lo que pone de manifiesto la relevancia 
estratégica del sector oleícola español tanto a nivel económico como tecnológico. Esta posición de liderazgo implica, 
además, una mayor exigencia en términos de control, trazabilidad y calidad del producto final.
(https://www.aceitedelasvaldesas.com/faq/varios/produccion-aceite-de-oliva-por-paises/?srsltid=AfmBOorTk2Hc_zO941GBcuBVAxz1h0AVcWQ8fFocKJPVzTscrF9gl97-)

El proceso de producción de la oliva no puede considerarse lineal ni estándar, ya que intervienen múltiples fases 
interconectadas y con un alto grado de variabilidad. Desde la gestión de parcelas agrícolas y la recolección de la 
aceituna, pasando por la recepción, el almacenamiento en depósitos o bodegas, la molturación, el envasado y finalmente 
la comercialización, cada etapa genera información crítica que debe ser registrada y relacionada con el resto del 
proceso. Asimismo, en este sector conviven distintas entidades, como agricultores, cooperativas, almazaras, envasadoras
y comercializadoras, cada una con funciones específicas pero dependientes entre sí, lo que incrementa la complejidad de 
la gestión de los datos.

A esta complejidad productiva se suma una regulación especialmente estricta, orientada a garantizar la seguridad 
alimentaria, la trazabilidad y la calidad del aceite de oliva. Normativas como el Reglamento CE 178/2002 obligan a 
disponer de sistemas que permitan identificar el origen del producto y reconstruir su recorrido en cualquier punto de 
la cadena. Además, los consejos reguladores, las denominaciones de origen, los criterios de clasificación del aceite 
según su categoría y los análisis físico-químicos y organolépticos, así como las certificaciones de calidad y producción 
ecológica, imponen requisitos adicionales que deben ser gestionados de forma precisa y documentada.

Las soluciones ERP generalistas presentan importantes limitaciones a la hora de adaptarse a este contexto, ya que no 
contemplan de forma nativa conceptos fundamentales del sector oleícola como la gestión por campañas agrícolas, los 
rendimientos grasos, los lotes dinámicos, las mezclas de calidades, los trasiegos entre depósitos o el control de mermas.
Esto provoca que su implantación requiera un elevado nivel de personalización, incrementando los costes, los plazos y el 
riesgo de fracaso del proyecto.

Como consecuencia de estas necesidades específicas, se hace imprescindible el desarrollo de software ERP diseñado 
específicamente para el sector de la producción de oliva. (Estos sistemas se caracterizan por estar orientados a modelos 
de gestión por campañas, permitir un control exhaustivo de la trazabilidad desde la parcela hasta el producto final, 
gestionar lotes de forma dinámica y automatizar la generación de informes oficiales exigidos por las administraciones y 
organismos reguladores. Además, facilitan el control de costes, la gestión de la calidad y la integración de todos los 
agentes implicados en la cadena productiva, permitiendo a cooperativas, agricultores y empresas maximizar el valor de su 
producción y mejorar su competitividad en el mercado. QUITAR ??) A continuación, se va a proceder a analizar los softwares
que realizan esta labor para ver cómo resuelven la problemática mencionada anteriormente.

\section{Análisis de Softwares ERP para la producción del aceite de oliva}

Actualmente, podemos encontrar diversos softwares ERP que están directamente relacionados con la producción de oliva,
otros que no se relacionan directamente pero se ubican dentro del sector agrícola y otros que son genéricos pero se 
pueden adaptar al sector. Vamos a proceder a comentar algunos de ellos:

\begin{itemize}
    \item \textbf{AgriCloud (RELACIÓN DIRECTA)}
    \newline
    AgriCloud es un Software de gestión agrícola diseñado para llevar un control de las fincas, trabajos, personal y
    facturación, pensada para agricultores, cooperativas, empresas de servicios agrícolas y trabajadores por campaña.
    AgriCloud destaca por ofrecer funcionalidades como las siguientes:
    \begin{itemize}
        \item Un cuaderno de campo digital según la normativa (muy importante como hemos mencionado anteriormente).
        \item Control horario y partes de trabajo desde el móvil.
        \item Facturación electrónica conectada con Hacienda, permitiendo llevar un control del personal y ver la
        rentabilidad real por finca o cultivo.
        \item Costes ingresos e informes automáticos en un solo lugar.
        \item Acceso desde cualquier dispositivo, tanto para móvil como tablet u ordenador.
        \item Generación de un cuaderno de campo digital, facilitando la gestión y el seguimiento de las actividades
        agrícolas de forma eficiente y precisa.
    \end{itemize} 

    Como hemos comentado, AgriCloud está pensada para poder satisfacer las necesidades de varios sectores:
    \begin{itemize}
        \item Agricultores individuales: podrán registrar en cuestión de segundos sus tareas pendientes o realizadas,
        llevarán un control de costes e ingresos por cultivo (sabiendo cuánto cuesta y cuánto produce por cada parcela)
        y podrán generar informes de forma automática.
        \item Cooperativas y asociaciones: destaca por aportar funcionales como el registro y consulta de la información
        de cada miembro y su producción, supervisión en tiempo real de cada tarea y de los resultados de la campaña,
        generación de informes globales o por socio en PDF o Excel, permite la comunicación y coordinación entre el 
        personal a través de un entorno digital único.
        \item Empresas de servicios agrícolas: las empresas podrán gestionar partes digitales por cliente o finca, 
        convertir dichos partes en facturas electrónicas al instante, control de maquinaria y materiales (horas y costes
        asociados a cada servicio o cliente), ayuda en la toma decisión gracias a la consulta de resultados por tipo de 
        trabajo y campañas.
        \item Trabajadores por campaña: AgriCloud ofrece una gestión de las "cuadrillas" de trabajadores a través de un 
        fichaje y control del horario digital, asignación por tareas por finca o cultivo, accesos personalizados (cada perfil
        está adaptado al rol del empleado) y ofrece un seguimiento en la nube, consiguiendo que toda la información se 
        sincronice automáticamente.
    \end{itemize}
    (enlace agricloud: https://www.agricloud.es/)

    \item \textbf{TuOlivar (RELACIÓN DIRECTA)}
    \newline
    TuOlivar es un Software de gestión agrícola para la producción del olivar. Se destacan principalmente por ofrecer las
    siguientes funcionalidades:
    \begin{itemize}
        \item Gestión de la explotación: permiten el control de fincas, parcelas, maquinaria, herramientas, mantenimiento
        de la maquinaria, trabajadores y campañas.
        \item Alta organización y uso desde cualquier ubicación: permite llevar un control de los trabajos realizados, control
        de gastos e ingresos y de la cosecha realizada, pudiendo agrupar la información por campañas, fincas y categorías.
        \item Gestión total: permite dar de alta a los clientes, almazaras, proveedores, notas, documentos o avisos, así como
        generar distintos informes para su control de costes y cosechas.
    \end{itemize}
    (enlace a tuolivar: https://www.tuolivar.com/)
    
    \item \textbf{Agroex}
    \newline
    Agroex es un Software de gestión agrícola industrial y ganadería, enfocado en empresas grandes. Cuentan con un software
    adaptado a las necesidades específicas del sector, permitiendo optimizar recursos y ahorrar costes (tanto en campaña
    agrícola específica como a lo largo del año), gracias a la administración y gestión tanto de los cultivos como del 
    personal. Además, cuenta con integraciones como AEMET o MAGRAMA, permitiendo el acceso a datos en tiempo real que 
    permite ayudar en la toma de decisiones.
    \newline
    Agroex ofrece soluciones para distintos sectores, desde un enfoque en la explotación agrícola y consultoría agraria
    a otros como bodegas, viñedos y viveros.
    \newline
    Si analizamos las funcionalidades que ofrece este Software, Agroex destaca por:
    \begin{itemize}
        \item Administración y gestión de explotaciones agrícolas y agroindustriales: permite la gestión de cultivos y 
        administración de explotaciones a través de cuadernos de campo, gestión de plazos, riegos y recursos hídricos.
        \item Conectividad con entidades como MAGRAMA o AEMET para obtener información meteorológica precisa y en tiempo real.
        \item Gestión de costes.
        \item Gestión contable, generando de manera automática los apuntes contables para cada movimiento realizado.
        \item Gestión de almacén, controlando el stock y gestión de mercancías.
        \item Gestión documental, incluyendo la gestión de maquinaria, productos fitosanitarios y gestión de personal.
        \item Gestión económica y tesorería, generando informes completos y exactos de la liquidez y resultados económicos.
        \item Gestión de producción, a través del control de las actividades de fabricación, producción y manufactura.
        \item Gestión eficiente de los recursos hídricos.
        \item Recursos humanos, gracias a una conexión con la Seguridad Social.
        \item Facturación, ofreciendo herramientas para gestionar presupuestos, seguimiento comercial y factura electrónica.
    \end{itemize}
    (https://agroex.es/)

    \item \textbf{Galdón Software (Enfoque en Cooperativas)}
    \newline
    Galdón Software es una empresa de ingeniería que ofrece soluciones software de gestión integrales para todo tipo de
    empresas e instituciones. Ofrecen servicios de consultoría y soluciones centralizadas para áreas de distintos
    negocios: ERP y CRM, Contabilidad y Finanzas, SGA, Apps, E-commerce y Recursos Humanos (RRHH).

    Dentro de los sectores en los que operan, desarrolladn Software ERP para almazaras y envasadoras de aceite. Con este
    ERP permiten llevar un control de la trazabilidad, producción y distribución del aceite y está pensado para cubrir
    las necesidades de cualquier empresa de aceite (oliva, girasol, soja, etc), almazara y envasadora.

    Entre sus características destacan:

    \begin{itemize}
        \item Planning de entrada y salida de mercancías.
        \item Control de necesidades de materias primas.
        \item Generación automática de partes de envasado y planificación asistida.
        \item Businnes Intelligence: cuentan con herramientas que permiten mostrar comparativas gráficas de ventas y 
        beneficios y análisis de ventas.
        \item Gesitón de cosecheros y puestos de compra.
    \end{itemize}

    Empresas reconocidas en el sector como Maeva (Granada) y Aceites Vallejo trabajan con este Software.
    (https://www.galdon.com/erp-almazaras-envasadoras-aceite/)

    \item \textbf{Campogest}
    \newline
    Campogest es un Software diseñado por la empresa Hispatec y destinado para los Técnicos Agrícolas, permitiéndoles 
    realizar una gestión de sus actividades a través del móvil o tablet.

    Esta aplicación destaca por ofrecer las siguientes funcionalidades:

    \begin{itemize}
        \item Gestión de fincas, cultivos, previsión meteorológica, tratamientos fitosanitarios y recomendaciones personalizadas
        a través del móvil.
        \item Generación de un cuaderno de campo teniendo en cuenta las recomendaciones en materia de Normativa Legal y 
        Protocolos de Calidad.
        \item Generación de planes de abonado y notificaciones a través del correo electrónico al agricultor.
        \item Posibilidad de generar la finca o parcela a través de la cartografía SIGPAC, pudiendo asignar un propietario.
        Gracias a esta funcionalidad los agricultores pueden dar de alta sus fincas y cultivos desde el mapa.
    \end{itemize}
    (https://www.campogest.com/)

    \item \textbf{Agroptima}
    \newline
    Agroptima es un Software de gestión agrícola cuyo objetivo es conseguir la simpleza a la hora de gestionar los cultivos
    y explotaciones por parte de agricultores y empresas agrícolas (según declara la propia empresa).

    En términos generales destaca por proporcionar un cuaderno de campo al agricultor, permitirle llevar un control de los 
    gastos y costes agrícolas (como los trabajadores) y saber en todo momento en qué estado se encuentra su finca.

    Entrando en más detalle, entre sus funcionalidades se encuentran:

    \begin{itemize}
        \item Localización de campos y cultivos, consulta de superficies y códigos SIGPAC.
        \item Anotación de toda la información que el agricultor considere relevante desde el móvil y sin cobertura.
        Entre esta información se destaca: trabajadores, maquinaria, horas trabajadas y productos.
        \item Toda actividad realizada y anotada quedará en el sistema y servirá para mejorar la toma de decisiones y 
        aumentar la rentabilidad de la gestión de la finca. 
        \item Las fincas o cultivos se podrán importar en la aplicación con un archivo Excel o Shapes, permitiendo visualizar
        las parcelas coloreadas según el tipo de cultivo y su superficie total.
        \item Gestión de roles, permitiendo que el usuario pueda gestionar aquellas actividades relacionadas con el.
    \end{itemize}

    Agroptima también destaca por ofrecer a sus clientes la posibilidad de generar un Cuaderno de Campo.

    (https://www.agroptima.com/)

    \item \textbf{Visual NACERT}
    \newline


    \item \textbf{AM System}
    \newline


    \item \textbf{BlueAiOlive}
    \newline 
    "BlueAiOlive es un software de gestión para almazaras y envasadoras de aceite, diseñado para cubrir todas
    las necesidades de gestión del sector agrícola". Así es como la empresa BlueAiCor define su ERP, destacando que a parte de cubrir
    las funcionalidades básicas de un ERP, tales como control y administración empresarial, también cuentan con módulos
    avanzados para la producción, envasado y comercialización del aceite de oliva.

    Para poder satisfacer dichas funcionalidades, cuentan con una herramienta integral que centraliza procesos, optimiza
    recursos y ofrece una visión global del negocio. De esta forma las empresas pueden tener un control total de cada etapa
    del proceso productivo.

    Entrando en el apartado del propio software, BlueAiOlive ofrece una app para oleicultores, donde se encontrará toda la
    información centralizada y disponible para el cliente a través de una interfaz sencilla e intuitiva. A través de ella,
    la empresa podrá llevar un control de la contabilidad, control del stock del almacém, facturación y liquidaciones, gestión
    de salidas de aceite, gestión de rendimiento de laboratorio (entre otras más).

    \newline \newline
    \textbf{FACTOR DIFERENCIAL - Conexión gestoría automatizada}
\end{itemize}



